% part: intuitionistic-logic
% chapter: semantics
% section: introduction

\documentclass[../../../include/open-logic-section]{subfiles}

\begin{document}

\olfileid[pt]{int}{sem}{int}

\olsection{Introdução}

Nenhuma lógica é descrita de modo satisfatório sem uma semântica, e a
lógica intuicionista não é exceção. Enquanto, para a lógica clássica,
a semântica baseada em !!{valuation}s é canônica, há várias semânticas
concorrentes para a lógica intuicionista. Nenhuma delas é completamente
satisfatória no sentido de oferecer uma explicação intuicionisticamente
aceitável dos significados dos conectivos.

A semântica baseada em !!{relational model}s, semelhante à semântica
das lógicas modais, talvez seja a mais popular. Nessa semântica,
!!{propositional variable}s são atribuídas a mundos, e esses mundos
são relacionados por uma relação de acessibilidade. Essa relação é
sempre uma ordem parcial, isto é, é reflexiva, antissimétrica e
transitiva.

Intuitivamente, podemos pensar nesses mundos como estados de
conhecimento ou ``situações evidenciais''. Um estado~$w'$ é acessível
a partir de~$w$ sse, tanto quanto sabemos, $w'$ é um possível estado
(futuro) de conhecimento, isto é, um estado compatível com o que se
sabe em~$w$. Uma vez conhecida, uma proposição não pode deixar de ser
conhecida: sempre que $!A$ é conhecida em~$w$ e~$Rww'$, $!A$ também é
conhecida em~$w'$. Assim, o ``conhecimento'' é monotônico em relação à
relação de acessibilidade.

Se definirmos ``$!A$ é conhecida'', como na lógica epistêmica, por
``verdadeira em todas as alternativas epistêmicas'', então
$!A \land !B$ é conhecida em~$w$ se tanto $!A$ quanto~$!B$ são
conhecidas em todas as alternativas epistêmicas. Mas, como o
conhecimento é monotônico e $R$ é reflexiva, isso significa que
$!A \land !B$ é conhecida em $w$ sse $!A$ e $!B$ são conhecidas
em~$w$. Pela mesma razão, $!A \lor !B$ é conhecida em $w$ sse pelo
menos uma delas é conhecida. Portanto, para $\land$ e $\lor$, as
condições de verdade dos conectivos coincidem com as da lógica clássica.

As condições de verdade do condicional, contudo, diferem das da lógica
clássica. $!A \lif !B$ é conhecida em~$w$ sse não existe $w'$ com
$Rww'$ no qual $!A$ seja conhecida sem que $!B$ também seja conhecida.
Isso não equivale à condição de que $!A$ seja desconhecida ou que
$!B$~seja conhecida em~$w$. Pois, se não conhecemos nem $!A$ nem $!B$
em~$w$, pode haver um estado epistêmico futuro~$w'$ com $Rww'$ tal que,
em $w'$, $!A$ seja conhecida sem que também passemos a conhecer~$!B$.

Conhecemos $\lnot !A$ apenas se não existe estado epistêmico futuro
possível no qual conheçamos~$!A$. A ideia é que, se $!A$ fosse
cognoscível, então $!A$ se tornaria conhecida em algum estado
epistêmico futuro possível. Como não podemos conhecer $\lfalse$, nesse
estado epistêmico futuro conheceríamos $!A$, mas não~$\lfalse$.

Nessa interpretação, o princípio do terceiro excluído falha. Pois há
alguns $!A$ que ainda não conhecemos, mas que podemos vir a conhecer.
Para tal !!a{formula}~$!A$, tanto $!A$ quanto~$\lnot !A$ são
desconhecidas; portanto, $!A \lor \lnot !A$ não é conhecida. Sabemos,
porém, por exemplo, que $\lnot(!A \land \lnot !A)$. Isso porque não é
possível um estado futuro no qual conheçamos tanto $!A$ quanto
$\lnot !A$, e sabemos disso independentemente de conhecermos ou não
$!A$ ou $\lnot !A$.

Os !!^{relational model}s não são a única semântica disponível para a
lógica intuicionista. A semântica topológica é outra: nela, as
proposições são interpretadas como conjuntos abertos em um espaço
topológico, e os conectivos são interpretados como operações sobre
esses conjuntos (por exemplo, $\land$ corresponde à interseção).

\end{document}
